% Exact alternative to termwise Abel summation; unresolved global implication.
\subsection*{Completing the radial frequency before modulus summation}

Let $q$ be coprime to $6$, let
$G_q=\sum_{x\bmod q}e(x^2/q)$, and write
\[
 T(m,n;q)=\sum_{a\bmod q}^{*}
       \left(\frac aq\right)e((ma+n\bar a)/q),\qquad
 S(m,n;q)=\sum_{a\bmod q}^{*}e((ma+n\bar a)/q).
\]
For arbitrary integers $d,h$ one has the exact finite Fourier identity
\begin{equation}\label{A:normal-completion}
 \sum_{k\bmod q}
   T(2,\overline{72}(k^2+d);q)e(-hk/q)
 =\left(\frac2q\right)G_q\,
       S(2-18h^2,\overline{72}d;q).
\end{equation}
Indeed, after interchanging the $k$ and $a$ sums, the inner
quadratic Gauss sum has coefficient
$A=\overline{72}\bar a$. Completing its square gives
\[
 \sum_{k\bmod q}e((Ak^2-hk)/q)
 =\left(\frac Aq\right)G_q e(-\overline{4A}h^2/q).
\]
Here $\overline{4A}=18a$ and
$(a/q)(A/q)=(2/q)$, proving
\eqref{A:normal-completion} for arbitrary odd composite $q$
as well as for primes.

For the decisive family $c=(1,0,k)$, the constant is $d=-18$,
so that $\overline{72}d=-\overline4$.
If $M=2-18h^2$, then
\begin{equation}\label{A:normal-classical}
 S(M,-\overline4;q)=-S(M,-1;2q).
\end{equation}
To check the sign and the modulus, Chinese remaindering gives
\[
 S(M,-1;2q)
 =S(M\bar q,-\bar q;2)\,
       S(M\overline2,-\overline2;q).
\]
The first factor is $(-1)^{M-1}=-1$ because $M$ is even.
In the second factor substitute the summation variable $2a$.
This gives \eqref{A:normal-classical}. Thus this family really
does reduce to ordinary integral-weight Kloosterman sums;
no extension of a named conjecture to a new multiplier is
being assumed for this identity.

Both identities are checked independently in exact cyclotomic
rings by \path{code/A_normal_completion_check.py}. The first
has 160 checks, including composite and prime-power moduli,
and the second has 40 checks. Their proofs above are valid
for every modulus stated.

\paragraph{Poisson completion with the endpoint phase retained.}
For $\psi$ Schwartz, $K>0$, and a real endpoint $X$, use
$\widehat\psi(\xi)=\int_{\mathbb R}\psi(v)e(-v\xi)\,dv$.
Poisson summation in residue classes gives exactly
\begin{align}
 &\sum_{k\in\mathbb Z}\psi(k/K)e(-Xk/q)
              T(2,\overline{72}(k^2+d);q)\nonumber\\
 &\quad=\frac Kq\left(\frac2q\right)G_q
     \sum_{h\in\mathbb Z}
       \widehat\psi\!\left(\frac{K(X-h)}q\right)
       S(2-18h^2,\overline{72}d;q).
 \label{A:normal-poisson}
\end{align}
The new index is centered at $h=X$, with width $q/K$.
This is a genuine improvement over discarding the endpoint
phase by absolute value: it turns the phase into the location
of the quadratic Kloosterman parameter.

For a dyadic endpoint-amplitude term
$K^{-1}\psi(k/K)e(-Xk/q)$, the factor $K$ on the right
of \eqref{A:normal-poisson} cancels.
The exact top-face denominator need not be expanded:
if $\eta$ is supported in a fixed positive dyadic interval, write
\[
 \frac{\eta(k/K)}{k+u}
   =K^{-1}\psi_{K,u}(k/K),\qquad
 \psi_{K,u}(v)=\frac{\eta(v)}{v+u/K}.
\]
For $K\geq1$ and $4\leq u\leq17/4$, this family and its Fourier
transforms have uniform Schwartz bounds. Thus the following
calculation applies to the exact top-face term, without silently
claiming a global bound for its separate radial remainder.
Multiplying the three original Gauss factors gives
$G_q^4=q^2$ and
$(6/q)(2/q)=(3/q)$. Consequently, after including the
$X^{-2}q^{-3}$ delta prefactor and the $X^3$ physical
volume factor, this dyadic term has the structural scale
\begin{equation}\label{A:normal-completed-scale}
 X\sum_h\sum_{q\asymp X}
    \frac{S(2-18h^2,-1;2q)}{q^2}
    \left(\frac3q\right)
    \widehat\psi\!\left(\frac{K(X-h)}q\right)
    A(q/X),
\end{equation}
up to fixed constants and the sign in
\eqref{A:normal-classical}. Here $A$ represents the fixed
face amplitude; its exact expression may be retained as
an integral over the fixed triangle.

In that face integral the remaining phase is $e(-Xu/q)$.
With $c=2q$, $m=18h^2-2$, and $n=1$, this has the standard
twisted Kloosterman form
\[
 e\!\left(\frac{2\sqrt{mn}\,\alpha}{c}\right),
 \qquad \alpha=-\frac{Xu}{\sqrt{18h^2-2}}.
\]
For $h\asymp X$ and $4<u<17/4$, $|\alpha|$ is bounded
independently of $X,K,h$ and is independent of $q$.
The Jacobi factor $(3/q)$ can be separated into finitely
many fixed residue classes. More explicitly, if $X\geq2$ and
$h\geq X/2$, then $|\alpha|\leq17/8$, so $B=3$ suffices.
The four good classes of $q$ modulo $12$ become the fixed
classes $c=2q\equiv2,10,14,22\pmod{24}$.
Conjecture 2 of Steiner \cite{A5Steiner2019} allows arbitrary
classes $c\equiv a\pmod{s}$, with no coprimality restriction
on $a,s$, and arbitrary nonzero indices. Its constants may
depend on $s,B$, which here are fixed at $24,3$.
This is why a bounded-twist
integral-weight conjecture is a more plausible input after
completion than it was before completion.

\paragraph{The remaining quantitative loss.}
Assume, only for this scaling test on the isolated top-face
family, Steiner's Conjecture 2 in the fixed classes just
specified. This does not assume a new multifrequency
extension of that conjecture.
Fix $\delta>0$ and consider $X^\delta\leq K\leq X^{1/2}$.
For $h$ outside $[X/2,2X]$, the Schwartz argument in
\eqref{A:normal-completed-scale} has size at least a
constant times $K$. Trivial bounds for the Kloosterman
sums, followed by an arbitrarily high Schwartz bound,
make this part smaller than any prescribed power of
$X$; choose the Schwartz order after $\delta$.
Thus the bounded-twist hypothesis is invoked only
where $h\asymp X$, and no uniformity for unbounded
$\alpha$ is being assumed.
For each $h$ in the effective range, Abel summation in
\eqref{A:normal-completed-scale} has weight $X/q$ of
bounded size and bounded variation. It therefore gives
$O_\epsilon(X^\epsilon)$ per $h$.
There are $O(X/K)$ effective $h$ values on $q\asymp X$,
with Schwartz tails treated dyadically. The resulting
termwise estimate is
\[
 O_\epsilon(X^{1+\epsilon}/K).
\]
This reaches square-root size for the terminal block
$K\asymp\sqrt X$, but does not do so uniformly for the
intermediate blocks.
The earlier hypothetical Sun estimate applied before
completion has scale $X^{1/2+\epsilon}K$ on the same
dyadic family. Taking the better of these two procedures
leaves
\[
 \min\{X^{1/2}K,\ X/K\},
\]
whose largest transition value is $X^{3/4}$ at
$K=X^{1/4}$.

This last comparison is an explicit accounting of two
termwise methods, not a lower bound for the true sum and
not an assertion that either conjecture fails to imply
the desired result. In particular, it does not complete
the requested known-conjecture implication. The next
concrete missing estimate is cancellation in the joint
$h,q$ sum in \eqref{A:normal-completed-scale}, with
$2-18h^2$ varying over a short quadratic sequence and
the bounded twist retained. A pointwise modulus bound
for each fixed $h$ supplies no such cancellation by
itself.
